% ============================================================
% 186_T0_Korrekturen_De_ch.tex
% Korrekturverzeichnis – FFGFT / T0-Zeit-Masse-Dualität
% Ältere Dokumente bleiben unverändert; dieses Dokument
% stellt Korrekturen und Präzisierungen klar.
% ============================================================

\begin{tcolorbox}[colback=red!3!white,colframe=red!50!black,boxrule=0.8pt,title={Korrekturhinweise},fonttitle=\bfseries]
\textbf{Korrekturen aus der Rechenprüfung (30.~Sept.~2026).} Die folgenden Stellen sind im Text korrigiert bzw. vermerkt; jede trägt dort einen datierten Hinweis mit dem vorherigen Wert:
\begin{itemize}
\item Gl.~(186.1) mit den angegebenen Eingaben: $\xi \approx 1{,}33\times 10^{-4} \to 1{,}304\times 10^{-4}$ ($-2{,}2\,\%$ gegenüber $4/30000$); Vorfaktor-Tabelle, Spalte „berechnet" mit $v = 246{,}22$~GeV: $0{,}511 / 105{,}7 / 1783 \to 0{,}505 / 105{,}1 / 1785$~MeV; $\hbar \sim \sqrt{\xi}\,\ell_P\,c$ in Gl.~(186.8) vermerkt (Dimension m²/s, zirkulär über $\ell_P$).
\item (1.~Okt.~2026) Herkunft des Vorfaktors $64\pi^4$: \glqq Tippfehler in der frühen HTML-Visualisierung\grqq{} $\to$ $64\pi^4 = 16\pi^3\cdot 4\pi$, das $4\pi$ aus $4\pi\varepsilon_0$ der Vakuumformel vom 6.~April 2025 ($= \xi_{\text{EFT}}\cdot\alpha$), beim Streichen von $\varepsilon_0$ nicht gekürzt (Dok.~385).
\end{itemize}
\end{tcolorbox}

\section*{Vorbemerkung}

Dieses Dokument listet Korrekturen und Präzisierungen zu bereits
veröffentlichten Dokumenten der FFGFT-Serie. Ältere Dokumente werden
\emph{nicht} verändert. Wo ein Widerspruch zwischen einem älteren und
einem neueren Dokument besteht, gilt die hier festgehaltene Version als
verbindlich.

Jeder Eintrag enthält: betroffenes Dokument und Stelle, den fehlerhaften
Ausdruck, den korrekten Ausdruck sowie eine kurze Begründung.

% ============================================================
\section{Korrektur~1: Vorfaktor in der Higgs-Ableitung von $\xi$}
% ============================================================

\subsection*{Betroffenes Dokument}

Dok.~049 (T0-Lagrangian und Higgs-Integration), HTML-Zwischenversion
sowie ältere Arbeitsversionen.

\subsection*{Fehlerhafter Ausdruck}

\begin{equation}
  \xi \;=\; \frac{\lambda_h^2\, v^2}{64\pi^4\, m_h^2}
  \;\approx\; 1{,}0 \times 10^{-5}
  \tag{alt -- falsch}
\end{equation}

\subsection*{Korrekter Ausdruck}

\begin{equation}
  \boxed{
    \xi \;=\; \frac{\lambda_h^2\, v^2}{16\pi^3\, m_h^2}
    \;\approx\; 1{,}30 \times 10^{-4}
  }
  \tag{186.1}
\end{equation}

mit den experimentellen Eingaben:
\begin{align*}
  m_h      &= 125{,}25\,\text{GeV} \quad \text{(Higgs-Masse)}, \\
  v        &= 246{,}22\,\text{GeV} \quad \text{(Vakuum-Erwartungswert)}, \\
  \lambda_h &= \frac{m_h^2}{2v^2} \approx 0{,}129
               \quad \text{(Higgs-Selbstkopplung)}.
\end{align*}

\subsection*{Begründung}

Der Vorfaktor $64\pi^4 = 16\pi^3 \cdot 4\pi$ stammt aus der Vakuumformel
vom 6.~April 2025,
$\lambda_h^2 v^2 e^2/(64\pi^4 \varepsilon_0 \hbar c\, m_h^2) = \xi_{\text{EFT}} \cdot \alpha$:
das $4\pi$ gehört zu $4\pi\varepsilon_0$ in $\alpha = e^2/(4\pi\varepsilon_0\hbar c)$ und wurde
beim Streichen von $\varepsilon_0$ (Übergang zu $\alpha = 1$) nicht mitgekürzt (Dok.~385).
\textit{(Korrigiert am 1.~Okt.~2026: vorher \glqq entstammt einem Tippfehler in der frühen HTML-Visualisierung\grqq{} -- $e^2/(4\pi\varepsilon_0\hbar c) = \alpha$, also $\lambda_h^2 v^2 e^2/(64\pi^4\varepsilon_0\hbar c\, m_h^2) = \lambda_h^2 v^2/(16\pi^3 m_h^2)\cdot\alpha$; ohne $\varepsilon_0$, $e$ und $\hbar c$ bleibt der überzählige Faktor $4\pi$ stehen.)}
Er liefert $\xi \approx 10^{-5}$, was um eine
Größenordnung vom geometrisch abgeleiteten Wert
$\xi = 4/3 \cdot 10^{-4}$ abweicht und die Naturkonstanten nicht
reproduziert. Der korrekte Vorfaktor $16\pi^3$ ergibt
$\xi \approx 1{,}30 \times 10^{-4}$, $2{,}2\,\%$ unter $4/30000$, in der Nähe des unabhängig
bestimmten Werts aus dem Proton-Elektron-Massenverhältnis (Dok.~009). \textit{(Korrigiert am 30.~Sept.~2026: vorher $\xi \approx 1{,}33 \times 10^{-4}$ in Gl.~(186.1) und hier -- mit $m_h = 125{,}25$~GeV, $v = 246{,}22$~GeV und $\lambda_h = m_h^2/(2v^2) = 0{,}1294$ ergibt $\lambda_h^2 v^2/(16\pi^3 m_h^2) = 1{,}304\times 10^{-4}$.)}

% ============================================================
\section{Korrektur~2: Tau-Yukawa-Vorfaktor $r_\tau$}
% ============================================================

\subsection*{Betroffenes Dokument}

Dok.~116 (Koide-Formel als Konsequenz der T0-Theorie), Tabelle der
Yukawa-Koeffizienten.

\subsection*{Fehlerhafter Ausdruck}

\begin{equation}
  r_\tau \;=\; \frac{8}{3} \;\approx\; 2{,}667
  \quad \Rightarrow \quad
  m_\tau \approx 1712\,\text{MeV} \quad (-3{,}6\,\% \text{ Abw.})
  \tag{alt -- falsch}
\end{equation}

\subsection*{Korrekter Ausdruck}

\begin{equation}
  \boxed{r_\tau \;=\; \frac{25}{9} \;\approx\; 2{,}778}
  \quad \Rightarrow \quad
  m_\tau \approx 1783\,\text{MeV} \quad (+0{,}4\,\% \text{ Abw.})
  \tag{186.2}
\end{equation}

Die Leptonmassen ergeben sich allgemein aus:
\begin{equation}
  m_\ell \;=\; r_\ell \cdot \xi^{p_\ell} \cdot v
  \tag{186.3}
\end{equation}
mit den Exponenten $p_e = 3/2$, $p_\mu = 1$, $p_\tau = 2/3$
(Schrittweite $\Delta p = 1/3$) und dem Vakuumerwartungswert $v$.

\subsection*{Begründung}

$r_\tau = 8/3$ war ein Überbleibsel aus einer früheren Iterationsstufe.
Dok.~158 (Koide-zu-g-2-Brücke) verwendet bereits $r_\tau = 25/9$, das
numerisch mit der Tau-Masse übereinstimmt. Zudem hat $25/9 = (5/3)^2$
eine geometrische Interpretation als Quadrat der reinen Sexte in der
Naturtonreihe, konsistent mit der musikalischen Resonanz-Analogie
(Dok.~060).

\subsubsection*{Vollständige korrigierte Vorfaktor-Tabelle}

\begin{center}
\begin{tabular}{lccc}
\hline
\textbf{Lepton} & \textbf{Exponent $p$} & \textbf{Vorfaktor $r$} & \textbf{$m$ [MeV], berechnet}\\
\hline
Elektron   & $3/2$  & $4/3$   & $0{,}505$ \\
Myon       & $1$    & $16/5$  & $105{,}1$ \\
Tau        & $2/3$  & $25/9$  & $1785$    \\
\hline
\end{tabular}
\end{center}
\textit{(Korrigiert am 30.~Sept.~2026: vorher $0{,}511$ / $105{,}7$ / $1783$~MeV -- $m = r\,\xi^p\,v$ mit $\xi = 4/30000$ und $v = 246{,}22$~GeV ergibt $0{,}5054$ / $105{,}05$ / $1785{,}0$~MeV; $0{,}511$ und $105{,}7$ sind die Messwerte, $-1{,}1\,\%$ bzw.\ $-0{,}6\,\%$ neben dem berechneten Wert.)}

% ============================================================
\section{Korrektur~3: Basisformel für $g$-$2$ des Elektrons}
% ============================================================

\subsection*{Betroffenes Dokument}

T0-Explorer (HTML-Visualisierung), Erstversion; ältere Zwischenversionen
von Dok.~018.

\subsection*{Fehlerhafter Ausdruck}

\begin{equation}
  a_e \;=\; \frac{4\pi}{f \cdot k_{\text{geom}}}
  \tag{alt -- falsch}
\end{equation}

\subsection*{Korrekter Ausdruck}

\begin{equation}
  \boxed{
    a_e \;=\; \frac{S_3}{f \cdot k_{\text{geom}}}
    \;=\; \frac{2\pi^2}{f \cdot k_{\text{geom}}}
  }
  \tag{186.4}
\end{equation}

mit $f = 1/\xi \approx 7500$, $k_{\text{geom}} = (2/\sqrt{\varphi})\cdot\sqrt{2} \approx 2{,}224$
und $S_3 = 2\pi^2 \approx 19{,}74$ (Oberfläche der 4D-Einheitskugel,
d.\,h.\ des 3D-Randes des Torus).

\subsection*{Begründung}

Der frühere Vorfaktor $4\pi$ entspricht der Oberfläche der 2D-Kugel im
3D-Raum, nicht der geometrisch korrekten 3D-Oberfläche des 4D-Torus.
$S_3 = 2\pi^2$ ist der korrekte Wert aus der Torus-Geometrie. Die
fraktalen Korrekturen für Myon und Tau verwenden weiterhin $4\pi$:

\begin{align}
  \Delta a_\mu  &= \frac{4\pi}{f^{5/3}}, \tag{186.5}\\
  \Delta a_\tau &= \frac{4\pi}{f^{4/3}}. \tag{186.6}
\end{align}

% ============================================================
\section{Präzisierung~1: Genauigkeit der Koide-Formel}
% ============================================================

\subsection*{Betroffene Dokumente}

Dok.~116 (Abstract, Z.\,14): $\Delta Q < 0{,}00003\,\%$ \\
Dok.~158 (Abstract): \glqq experimentell bestätigt auf $0{,}001\,\%$\grqq

\subsection*{Korrekte Einordnung}

Die Koide-Relation
\begin{equation}
  Q \;=\; \frac{m_e + m_\mu + m_\tau}
               {(\sqrt{m_e}+\sqrt{m_\mu}+\sqrt{m_\tau})^2}
  \;=\; \frac{2}{3}
  \tag{186.7}
\end{equation}
ist experimentell auf besser als $0{,}001\,\%$ bestätigt
(PDG-Werte 2022). Die Angabe $< 0{,}00003\,\%$ in Dok.~116 bezieht
sich auf die interne Konsistenz der T0-Formeln, nicht auf die
experimentelle Messunsicherheit. Beide Aussagen sind inhaltlich korrekt,
aber unterschiedliche Größen. In externen Kommunikationen ist
$0{,}001\,\%$ die geeignete Angabe.

\subsection*{Hinweis}

Die T0-Einzelmassen stimmen auf $\sim 1\,\%$ mit dem Experiment überein;
die hohe Präzision von $Q = 2/3$ entsteht durch die spezifische
algebraische Struktur der Koide-Kombination, nicht durch individuelle
Massengenauigkeit.

% ============================================================
\section{Präzisierung~2: $h$ aus dem Higgs-Feld -- Ableitungskette}
% ============================================================

\subsection*{Kontext}

Die Planck-Konstante $\hbar$ wird in verschiedenen Dokumenten
unterschiedlich hergeleitet. Hier wird die vollständige, konsistente
Ableitungskette festgehalten.

\subsection*{Verbindliche Ableitungskette}

\begin{equation}
  \xi \;=\; \frac{\lambda_h^2 v^2}{16\pi^3 m_h^2}
  \;\longrightarrow\;
  L_0 \;=\; \xi \cdot \ell_P
  \;\longrightarrow\;
  \hbar \;\sim\; \sqrt{\xi}\cdot \ell_P \cdot c
  \tag{186.8}
\end{equation}
\textit{(Vermerk vom 30.~Sept.~2026: $\sqrt{\xi}\cdot\ell_P\cdot c$ hat die Einheit m²/s, nicht J\,s; numerisch $0{,}01155 \cdot 1{,}616\times 10^{-35}\,\text{m} \cdot 3\times 10^{8}\,\text{m/s} = 5{,}6\times 10^{-29}$~m²/s. Zudem enthält $\ell_P = \sqrt{\hbar G/c^3}$ selbst $\hbar$; die Kette ist in dieser Form zirkulär.)}

Die Zeit-Masse-Bindungsbedingung liefert die physikalische Begründung:
\begin{equation}
  T(x,t)\cdot m(x,t) \;=\; 1
  \quad\Longrightarrow\quad
  \Tfield \;=\; \frac{\hbar}{m\, c^2}
  \quad\Longrightarrow\quad
  E \;=\; \hbar\omega.
  \tag{186.9}
\end{equation}

$h = 2\pi\hbar$ ist damit \emph{keine} freie Konstante, sondern eine
geometrische Konsequenz aus $\xi$ und $\ell_P$. Der numerische SI-Wert
$h \approx 6{,}626 \times 10^{-34}\,\text{J·s}$ ist das
Kalibrierungsergebnis dieser Beziehung auf die SI-Basiseinheiten.

% ============================================================
\section*{Zusammenfassung der Korrekturen}
% ============================================================

\begin{center}
\begin{tabular}{p{0.08\linewidth}p{0.22\linewidth}p{0.28\linewidth}p{0.28\linewidth}}
\hline
\textbf{Nr.} & \textbf{Dok.} & \textbf{Fehlerhaft} & \textbf{Korrekt} \\
\hline
K1 & 049 (HTML) & $\xi = \lambda_h^2 v^2/(64\pi^4 m_h^2)$ &
     $\xi = \lambda_h^2 v^2/(16\pi^3 m_h^2)$ \\[4pt]
K2 & 116 (Tab.) & $r_\tau = 8/3$ &
     $r_\tau = 25/9$ \\[4pt]
K3 & 018 (Explorer) & $a_e = 4\pi/(f\cdot k)$ &
     $a_e = 2\pi^2/(f\cdot k)$ \\[4pt]
P1 & 116 / 158 & $\Delta Q < 0{,}00003\,\%$ vs.\ $0{,}001\,\%$ &
     Beide korrekt, verschiedene Größen; $0{,}001\,\%$ für extern \\[4pt]
P2 & mehrere & $\hbar$ als Postulat &
     $\hbar$ folgt aus $\xi$ via $L_0 = \xi\cdot\ell_P$ \\
\hline
\end{tabular}
\end{center}
